
\documentclass[12pt, letterpaper, twoside]{article}
\usepackage[utf8]{inputenc}
\usepackage{enumerate} % enumerados

\usepackage{float} % para usar [H]

\title{Organizaci\'on y Arquitectura de computadoras Práctica 1}
\author{José Fernando Reyes García}
\date{18 de febrero de 2017}
 
\begin{document}

\maketitle

Nombre del alumno: José Fernando Reyes García
Datos de la computadora:
\begin{itemize}
\item Fabricante y modelo de la computadora:
  
  Hewlett Packard, HP Notebook .
  
\item Fabricante, modelo, frecuencia, número de núcleos, y arquitectura del procesador:
  
  AMD A8-7410 APU, 2.20GHz, 4 núcleos, 64 bits.
  
\item Capacidad de memoria RAM y de cachés de los procesadores:
  
  4 GB Memoria RAM, 2 MB memoria caché del procesador.
    
\item Capacidad del disco duro:
  
  1 TB de capacidad.
  
\item Distribución de linux y versión del kernel:
  
  Ubuntu 16.04, Kernel: 4.4.0-57-generic. 
\end{itemize}

\begin{table}[H]
\begin{center}
\begin{tabular}{|l|l|}
\hline
Nombre de la prueba & Resultado \\
\hline \hline
GZIP Compression & 30.56 Seconds \\ \hline
DCRAW & 145.45 Seconds \\ \hline
FLAC Audio Encoding & 21.75 Seconds \\ \hline
GnuPG & 34.88 Seconds\\  \hline
REDIS & 385728.07 Requests Per Second \\ \hline
Timed MAFFT Alignment & 19.45 Seconds \\ \hline
Timed MrBayes Analysis & 46.58 Seconds \\ \hline
Timed M Player Compilation & 263.09 Seconds \\ \hline
Timed PHP Compilation & 128.88 Seconds\\ \hline
\end{tabular}
\caption{Resultado de las pruebas}
\label{tabla:sencilla}
\end{center}
\end{table}

\end{document}

